\chapter{{\fontsize{14pt}{21pt}\selectfont\MakeUppercase{Аналоги и конкуренты}}}
\label{ch:analogues}

\section{Выявленные аналоги и конкуренты}

Рассматриваются аналоги, которые частично закрывают контур управления цифровыми каналами МСБ. Перечень охватывает три основных класса решений: CRM и омниканальные коммуникации (Bitrix24, SendPulse, Kommo), SMM-планирование и публикации (Hootsuite, Buffer, SMMplanner), мониторинг упоминаний и ORM (YouScan).

Ни одно из рассмотренных решений не обеспечивает поддержания <<эталонных>> данных о~бизнесе и контроля расхождений между каналами как отдельной функции, единого протокола оркестрации для платформ с~различным уровнем открытости интерфейсов (API и RPA), интеграции с~ключевыми для российского рынка геосервисами (Яндекс.Бизнес, 2ГИС), не имеющими публичных API, а~также сквозного журналирования результатов по всем каналам.

\section{Критерии для сравнения аналогов и конкурентов}

Для сравнения используется метод взвешенных критериев: каждому критерию присваивается удельный вес, каждый аналог оценивается по шкале от~0 до~3 баллов (0~--- функция отсутствует, 1~--- минимально, 2~--- частично, 3~--- полноценно). Итоговый балл~--- взвешенная сумма.

\begin{itemize}
  \item \textbf{К1. Функциональный охват} (вес 0{,}20): CRM и контакт-центр; планирование публикаций; мониторинг упоминаний и репутации.
  \item \textbf{К2. Мультиканальность} (вес 0{,}15): количество поддерживаемых каналов и степень <<единого окна>>.
  \item \textbf{К3. Автоматизация} (вес 0{,}15): встроенные сценарии, правила, массовые операции.
  \item \textbf{К4. Управление консистентностью данных} (вес 0{,}25): поддержка <<эталона>> и контроль расхождений~--- ключевой критерий.
  \item \textbf{К5. Трассируемость} (вес 0{,}10): журналирование действий и статусов.
  \item \textbf{К6. Расширяемость} (вес 0{,}15): API и интеграции, возможность добавления каналов, в~т.\,ч. для платформ без API.
\end{itemize}

Результаты балльной оценки приведены в~таблице~\ref{tab:analogues}. Оценки OneVoice являются проектными.

\begin{table}[H]
  \caption{Балльная оценка аналогов по критериям сравнения}
  \label{tab:analogues}
  \small
  \begin{tabularx}{\textwidth}{|X|c|c|c|c|c|c|c|}
    \hline
    \textbf{Решение} & \textbf{К1} & \textbf{К2} & \textbf{К3} & \textbf{К4} & \textbf{К5} & \textbf{К6} & \textbf{Итог} \\
    \textbf{} & \textbf{0{,}20} & \textbf{0{,}15} & \textbf{0{,}15} & \textbf{0{,}25} & \textbf{0{,}10} & \textbf{0{,}15} & \\
    \hline
    Bitrix24 & 3 & 2 & 2 & 1 & 2 & 2 & 1{,}95 \\
    \hline
    SendPulse & 2 & 2 & 2 & 0 & 1 & 2 & 1{,}40 \\
    \hline
    Kommo & 2 & 2 & 2 & 0 & 1 & 2 & 1{,}40 \\
    \hline
    Hootsuite & 1 & 2 & 2 & 0 & 1 & 2 & 1{,}20 \\
    \hline
    Buffer & 1 & 2 & 2 & 0 & 1 & 2 & 1{,}20 \\
    \hline
    SMMplanner & 1 & 1 & 2 & 0 & 0 & 1 & 0{,}80 \\
    \hline
    YouScan & 1 & 1 & 1 & 0 & 1 & 2 & 0{,}90 \\
    \hline
    \textbf{OneVoice*} & \textbf{2} & \textbf{2} & \textbf{3} & \textbf{3} & \textbf{2} & \textbf{3} & \textbf{2{,}55} \\
    \hline
  \end{tabularx}
\end{table}

\noindent *~Проектные оценки OneVoice. По ключевому критерию К4 (консистентность данных) все аналоги получили 0--1 балл; OneVoice проектируется с~эталонным профилем и автоматическим обнаружением расхождений, а~по К6~--- с~гибридной моделью интеграции (API + RPA для платформ без API).
